
\section{The original line-to-cusp argument}
\label{sec:sr-source-line-cusp}

This section translates
\srcid{EigenvalueEndgame.CuspidalAverages.CuspOfLineOfKeys},
together with the coefficient lemma \srcid{KillCoef.kill\_coef}.
It retains the two source assertions \(\operatorname{KeyTwo}\)
and \(\operatorname{KeyThree}\). It does not replace them by
a general-linear Jacquet-module or unitarity argument.

For a rational isotropic \(j\)-frame \(u\) with dual frame \(u'\)
in \(W\), write \(U=\operatorname{colSp}(u)\) and let
\(\operatorname{ct}_U\) be the normalized average along the
source group \(\operatorname{unipU}(W,U)\).
The group fixes \(U\) pointwise and acts trivially on \(U^\perp/U\).
Its rational quotient is compact.
The context \(\operatorname{Ctx}\) used below comprises the CM
involution, skew symmetry of the target form, and its nonzero
determinant. From these the source constructs the quotient
averages and rational frame extensions.

A continuous additive character \(\chi\) of \(\mathbb A_F\otimes_F E\)
has the source property \(\operatorname{IsChar}\) if, for every
split finite place and each selected projection
\(t:E\to F_v\), it is nontrivial on some vector supported at \(v\)
whose conjugate projection is zero.
A functional \(\Lambda\) has the source property
\(\operatorname{IsWh}(W,u,\chi)\) if it is additive on continuous
rationally invariant functions, is homogeneous, and obeys
\[
 \Lambda(R_hf)=
       \chi\!\left(\sum_{i=0}^{j-2}g_{i,i+1}\right)\Lambda(f)
 \tag{C1}\label{eq:sr-source-whittaker-law}
\]
whenever \(h\) acts on \(u\) by an upper unipotent matrix \(g\)
and acts trivially on \(U^\perp/U\).
Here \(R_hf(x)=f(xh)\).
Both the condition on the complementary quotient and the order
of right translation in (C1) are part of the assertion.

The retained \(\operatorname{KeyThree}\) says that for every
allowable triple, every such frame of size \(j\ge3\),
every admissible \(\chi,\Lambda\), and every
\(f\in\operatorname{Sig}(G)\),
\[
                         \Lambda(R_hf)=0
                              \quad\hbox{for every }h.
 \tag{C2}\label{eq:sr-source-keyThree}
\]
The retained \(\operatorname{KeyTwo}\) has a different conclusion.
For a two-frame, its specified continuous Levi map \(m\),
admissible \(\chi,\Lambda\), \(f\in\operatorname{Sig}(G)\),
and any \(h_0\), there exist a neighborhood \(V\) of \(h_0\),
matrices \(b_i\in\mathrm{GL}_2(\mathbb A_F\otimes_F E)\),
and numbers \(y_i\in\mathbb C\) such that
\[
 \begin{gathered}
 b_0=1,\qquad \det b_i=1,\qquad |y_i|\longrightarrow\infty,\\
 \forall h\in V\ \exists c_h\in\mathbb C\
       \forall i\ge0,\qquad
      \Lambda(R_{m(b_i)h}f)=c_h y_i .
 \end{gathered}                                    \tag{C3}
 \label{eq:sr-source-keyTwo}
\]
The same \(b_i,y_i\) work throughout \(V\); \(c_h\) may vary with \(h\).
This neighborhood quantifier is used in convolution below.

\begin{srlemma}[The coefficient version of the growth bootstrap]
\label{lem:sr-source-kill-coefficient}
Use the abstract \(G,X,\Gamma,n,a,\lambda\), cusp condition,
and growth notion of
Lemma~\ref{lem:sr-source-abstract-contraction}.
Let \(\phi,c\) satisfy its data with recurrence parameter \(z=0\),
and let \(\phi\) have some finite growth exponent.
Let \(e:X\to\mathbb C\) be continuous with \(|e(x)|\le1\).
Suppose \(|y_i|\to\infty\), \(b_i\in G\), and on a neighborhood
\(V\) of \(1\),
\[
 y_0\int_Xe(x)\phi(n(x)b_i z)\,d\lambda(x)
       =y_i\int_Xe(x)\phi(n(x)z)\,d\lambda(x)
                    \qquad(z\in V,\ i\ge0).
 \tag{C4}\label{eq:sr-source-weighted-coefficient-scaling}
\]
Then \(\int_Xe(x)\phi(n(x))\,d\lambda(x)=0\).
\end{srlemma}

\begin{proof}
This is \srcid{KillCoef.kill\_coef}, using the same
growth induction \srcid{ConvCont.kill\_ind}.
If \(\phi\) is globally bounded by \(M\), every weighted
probability average has absolute value at most \(M\).
If its value at \(1\) were nonzero, (C4) at \(z=1\)
would contradict \(|y_i|\to\infty\).

For the induction step set
\(w(g)=\int_Xe(x)\phi(n(x)g)\,d\lambda(x)\).
The compact-set bound \srcid{Dat.bound} and dominated
convergence make \(w\) continuous.
Choose \(V'\ni1\) with \(V'V'\subset V\).
Take the growth-lowering bump with support in an arbitrary
prescribed neighborhood intersected with \(V'\).
Fubini, with the compact-set bound times the bump as
majorant, gives
\[
       \int_Xe(x)\phi_\eta(n(x)g)\,d\lambda(x)
                                  =(w*\eta)(g).
\]
For \(z\in V'\) and \(u\) in the bump support, \(zu\in V\).
Integrating (C4) at \(zu\) shows that the convolved
function has the same scaling relation on \(V'\).
The lower-growth induction makes its weighted coefficient
at \(1\) zero. The positive-bump argument
\srcid{Bump.zero}, applied to the continuous \(w\),
then gives \(w(1)=0\).
No spectral decomposition is involved.
\end{proof}

Fix a standard adelic additive character \(\psi\) of \(E\).
The context's adelic identification
\(\mathrm{ad}:\mathbb A_F\otimes_F E\to\mathbb A_E\)
gives \(\chi=\psi\circ\mathrm{ad}\).
The source proves \(\operatorname{IsChar}(\chi)\):
if a split component were annihilated, use the continuous
split section with projections \((1,0)\).
The standard character's rational-annihilator property would
force that adelic section to be rational, contradicting
its two projections. This is \srcid{isChar}.

For the chosen \(j\)-frame define \(\Lambda_U(f)\) as the full
Whittaker coefficient, evaluated at \(1\), of
\[
       g\longmapsto
          \operatorname{ct}_U(f)\bigl(\operatorname{lev}_u
                                      (\mathrm{ad}^{-1}g)\bigr)
                \quad\hbox{on }\mathrm{GL}_j(\mathbb A_E).
 \tag{C5}\label{eq:sr-source-whittaker-of-average}
\]
In the source this is
\(\operatorname{whit}_{\psi,j-1}\) and is named \srcid{wh}.
Its properties, with their precise uses, are as follows.

First, \srcid{isWh} verifies (C1) using the upper-unipotent
Whittaker transformation law. The constant-term average
supplies the needed complementary-quotient invariance.
Second, \srcid{wh\_right} identifies
\(\Lambda_U(R_{\operatorname{lev}_u(g)h}f)\) with the value
at \(g\) of the Whittaker coefficient of the \(h\)-translated
Levi restriction.
Third, \srcid{wh\_ne} is the detection implication:
if averages along all proper nonzero initial subframes
of \(u\) vanish, and \(\Lambda_U(R_hf)=0\) for every \(h\),
then \(\operatorname{ct}_U(f)=0\).
Indeed \srcid{ct\_van} identifies each proper general-linear
constant term of the Levi restriction with one of those
smaller averages. The source theorem \srcid{whit\_good}
then detects a nonzero Levi restriction by its full Whittaker
coefficient. This use of the Whittaker expansion does not
introduce a cuspidal-spectrum unitarity argument.

\begin{srlemma}[The two-frame coefficient, \srcid{wh\_two}]
\label{lem:sr-source-wh-two}
Let \(f\) be continuous, rationally invariant and of polynomial
adelic height growth. Suppose its average along the first
line of a rational two-frame vanishes.
If its Whittaker functional has the neighborhood scaling
property (C3), then \(\Lambda_U(R_{h_0}f)=0\).
\end{srlemma}

\begin{proof}
Embed \(\mathrm{SL}_2(\mathbb A_E)\) in the Levi of the two-frame.
The matrices \(b_i\), having determinant one, lie in this group.
On it use
\(\phi(g)=\operatorname{ct}_U(f)(m(g)h_0)\).
Its upper-unipotent mean is zero by the vanishing first-line
average and integration in stages.
Thus it satisfies the abstract data with \(c=0,z=0\).
It has finite ray growth because the original \(f\) has
polynomial height growth and the two-plane average uses a
compact rational quotient.
The concrete cusp/bump condition for this \(\mathrm{SL}_2\)
comes from \srcid{AeFacts.ae\_facts}.

Write the Whittaker coefficient as its compact additive
quotient integral with weight
\[
                 e(x)=\psi(-\operatorname{dot}(e_{\rm last},x)).
\]
This weight is continuous and has absolute value one.

Transport the neighborhood \(V\) in (C3) to a neighborhood
of \(1\) under \(g\mapsto m(g)h_0\).
The relation \(b_0=1\) gives (C4), by multiplying the
\(i\)-th identity by \(y_0\) and the zeroth identity
by \(y_i\). The preceding coefficient lemma makes the
Whittaker coefficient zero at \(h_0\), as asserted.
\end{proof}

\begin{srproposition}[The source theorem \srcid{cuspOfLine\_of\_keys}]
\label{prop:sr-source-line-cusp}
Assume \(\operatorname{WeilConv}\),
\(\operatorname{KDecay}\), \(\operatorname{IAut}\),
\(\operatorname{KeyThree}\), and \(\operatorname{KeyTwo}\).
For an allowable triple, let \(H'\) be any nondegenerate
rational target plane and \(f\in\operatorname{Sig}(G)\).
If its line constant term at \(H'\) is zero, then \(f\)
is cuspidal in the source's sense: every average along
\(\operatorname{unipU}(W,U)\), for a nonzero rational
isotropic \(U\), vanishes for every Haar measure and
rational fundamental domain.
\end{srproposition}

\begin{proof}
Automorphy follows from \(\operatorname{IAut}\).
The source lemma \srcid{sig\_mod} combines the large
\(\operatorname{WeilConv}\) bound and the small
\(\operatorname{KDecay}\) bound into polynomial growth
of their difference. Its height is at least one, so the
two powers are bounded by their summed exponent after
enlarging their nonnegative constants.

We strongly induct on the positive frame dimension \(j\),
proving zero normalized average for every rational
isotropic \(j\)-frame with a dual frame.
For \(j=1\), \srcid{AvLine.av\_line} transports the given
zero line average to every other rational isotropic line.
Witt extension supplies a rational isometry between their
nonzero vectors, \srcid{conj\_N} identifies the two
unipotent groups, and automorphy and normalized quotient
integration transport the zero assertion.

For \(j=2\), the induction hypothesis gives the zero first-line
average. Apply \(\operatorname{KeyTwo}\) to the functional
(C5) and to each \(h_0\).
Lemma~\ref{lem:sr-source-wh-two} makes all its right translates
zero. The detection implication \srcid{wh\_ne} therefore
gives the zero two-plane average.

For \(j\ge3\), the proper initial-subframe averages vanish by
induction. The functional (C5) satisfies
\(\operatorname{IsWh}\), and \(\operatorname{KeyThree}\)
makes every one of its right translates zero.
Apply \srcid{wh\_ne} again to obtain the zero \(j\)-plane
average.

Finally any nonzero rational isotropic subspace \(U\)
has a basis frame and an isotropic dual frame:
\srcid{Ctx.frame} obtains the dual from nondegeneracy.
The preceding induction proves its normalized average is zero.
To obtain the exact unnormalized-integral quantifiers,
\srcid{Ctx.int\_zero} scales an arbitrary Haar measure by the
inverse mass of its fundamental domain. The rational quotient
is compact, so this mass is finite and positive.
Canceling the resulting nonzero scalar proves zero for
the originally given measure and domain.
\end{proof}

The inputs \(\operatorname{KeyTwo}\) and
\(\operatorname{KeyThree}\) are produced separately in
\srcid{CuspidalAverages.ActsOfReads}:
\srcid{keyThree\_of\_locThree} uses
\(\operatorname{ThetaBasic},\operatorname{KDecay},
\operatorname{IAut},\operatorname{LocThree},\operatorname{Places}\);
\srcid{keyTwo\_of\_locTwo\_vacBar} uses the corresponding
\(\operatorname{LocTwo},\operatorname{VacBar},\operatorname{Places}\)
inputs. Their local support, torus-scaling, vacuum, and
split-place producers have not been silently replaced by
a different proof in this section.

\subsection{The source's local support and three-column distribution argument}
\label{subsec:sr-source-local-three}

This subsection gives the actual mechanism of
\srcid{ZariskiLocallyAffineLine} and
\srcid{LocalFieldT2Space.locSupp/locThree}.
It uses algebraic linear functionals on locally constant compactly
supported functions. No continuity, distribution completion, or
unitary representation assumption on the functional is made.

Let \(k\) be a nonarchimedean local field, \(\mathcal O\) its valuation
ring, and \(V=M_{n,d}(k)\).
The subgroup \(\Lambda_0=M_{n,d}(\mathcal O)\) is compact and open:
openness follows entry by entry from the open valuation ring, and
compactness from the finite product of compact closed unit balls.
The field is Hausdorff, and every neighborhood of zero contains a
nonzero power of a uniformizer.
Write \(\mathcal S(V)\) for locally constant compactly supported
complex functions.
Let \(l\) be a complex-linear functional on all functions on \(V\);
only its restriction to \(\mathcal S(V)\) matters.
Say that \(l\) kills a subset \(U\) if \(l(\varphi)=0\) for every
\(\varphi\in\mathcal S(V)\) whose nonzero set lies in \(U\).

\begin{srlemma}[Compact cosets and local vanishing]
\label{lem:sr-source-local-kills}
If every \(Y\in U\) has a neighborhood killed by \(l\), then \(l\)
kills \(U\).
More generally let \(\Lambda\subset V\) be a compact open additive
subgroup, and suppose \(l\) kills \(U\).
Suppose that for every \(Y\) in a neighborhood \(\Omega\), either
\(Y\in U\), or there are a \(k\)-linear map \(s_Y:V\to V\)
and \(c_Y\ne1\) with
\[
 s_YY=Y,\qquad
 s_YZ\in\Lambda\ \Longleftrightarrow\ Z\in\Lambda,\qquad
 l(\varphi\circ s_Y)=c_Yl(\varphi)\quad(\varphi\in\mathcal S(V)).
 \tag{LD1}\label{eq:sr-source-local-stabilizer}
\]
Then \(l\) kills the interior of \(\Omega\).
\end{srlemma}
\begin{proof}
A locally constant compactly supported function is a finite sum
of constants times indicators of cosets of a sufficiently small
compact open subgroup \(a\Lambda\), \(a\in k^\times\).
Here is the precise smallness argument.
Its nonzero set is compact: it is closed, because the zero fibre of
a locally constant function is open, and is contained in its compact
support.
At each point of this set, local constancy supplies a neighborhood
on which its value is constant.
Compactness and continuity of addition and scalar multiplication
give one sufficiently small \(a\ne0\) such that every coset
\(x+a\Lambda\) meeting this set is contained in an assigned such
neighborhood.
For a function supported in an open set, impose containment in that
open set at the same time.
The compact nonzero set meets only finitely many disjoint open
cosets. Summing its constant values on those cosets proves the
decomposition. This also proves the first assertion by choosing
the assigned neighborhoods to be killed by \(l\).

For the second assertion take \(\varphi\) whose nonzero set is
contained in \(\operatorname{int}\Omega\), and choose such a coset
decomposition with every relevant coset contained in \(\Omega\).
For a coset \(C=x+a\Lambda\), either \(C\subset U\), in which case
its indicator has zero \(l\)-value, or choose \(Y\in C\setminus U\).
Since \(s_Y\) is linear, fixes \(Y\), and preserves membership in
\(\Lambda\) in both directions, it preserves membership in
\(Y+a\Lambda=C\) in both directions.
Thus \(1_C\circ s_Y=1_C\).
Equation (LD1) gives
\((c_Y-1)l(1_C)=0\), so this value is zero.
Add the finitely many coset values.
These are \srcid{small}, \srcid{kills\_of\_local}, and
\srcid{kills\_interior} in the source.
\end{proof}

\begin{srlemma}[The source's moving square-zero stabilizers]
\label{lem:sr-source-local-dead}
Suppose \(l\) kills \(U\).
Fix \(Y_0\in V\).
Let \(A_Y,B_Y:V\to V\) be linear maps with
\((Y,Z)\mapsto A_YZ,B_YZ\) continuous at every \((Y_0,Z)\).
Assume
\[
 A_{Y_0}^2=B_{Y_0}^2=0,\qquad
                       A_{Y_0}B_{Y_0}=B_{Y_0}A_{Y_0}.
\]
For every \(Y\) sufficiently near \(Y_0\) outside \(U\), suppose
\[
 \begin{gathered}
 A_Y^2=B_Y^2=0,\qquad
            (1+A_Y)(1+B_Y)Y=Y,\\
 \exists c_Y\ne1\quad
 l\bigl(\varphi\circ(1+A_Y)(1+B_Y)\bigr)
                  =c_Yl(\varphi)\quad(\varphi\in\mathcal S(V)).
 \end{gathered}\tag{LD2}\label{eq:sr-source-local-moving-stabilizer}
\]
Then \(l\) kills a neighborhood of \(Y_0\).
\end{srlemma}
\begin{proof}
Set
\[
 \Lambda=\Lambda_0\cap A_{Y_0}^{-1}\Lambda_0
       \cap B_{Y_0}^{-1}\Lambda_0
       \cap(A_{Y_0}B_{Y_0})^{-1}\Lambda_0.
\]
It is open and closed, and is a closed subset of the compact
\(\Lambda_0\), hence compact.
The square-zero and commutation equations show that both
\(A_{Y_0}\) and \(B_{Y_0}\) send \(\Lambda\) into itself.
For example the four conditions for \(A_{Y_0}Z\) are supplied by
\(A_{Y_0}Z,0,A_{Y_0}B_{Y_0}Z,0\in\Lambda_0\).

Compactness of \(\Lambda\), openness of \(\Lambda\), and the stated
joint continuity show that, after shrinking around \(Y_0\),
\[
                       A_Y\Lambda,\ B_Y\Lambda\subset\Lambda.
\]
Explicitly, at each \(Z\in\Lambda\), continuity gives product
neighborhoods of \(Y_0,Z\) whose images lie in \(\Lambda\);
a finite cover of the compact set of \(Z\)'s gives one
neighborhood of \(Y_0\) valid for every \(Z\in\Lambda\).
If \(C\Lambda\subset\Lambda\) and \(C^2=0\), then
\[
                     Z+CZ\in\Lambda\ \Longleftrightarrow\
                                     Z\in\Lambda.
\]
For the converse subtract \(C(Z+CZ)=CZ\).
Apply this to \(C=A_Y,B_Y\).
Thus \(s_Y=(1+A_Y)(1+B_Y)\) satisfies all clauses of (LD1)
outside \(U\). The preceding lemma, on the interior of the chosen
neighborhood, proves the result.
This is the actual argument of \srcid{dead} and \srcid{deadM}.
\end{proof}

Fix distinct indices \(i_0,i_1\).
For a row \(\rho\in k^d\), put
\[
 C_{i,\rho}(Z)=Z_{\cdot,i}\rho,\qquad
                     \mathrm{colOp}_{i,\rho}Z=Z+C_{i,\rho}(Z).
\]
If \(\rho'_i=0\), then
\(C_{i,\rho}C_{j,\rho'}=0\), by multiplication of the two
rank-one matrices.
In particular \(C_{i,\rho}^2=0\) when \(\rho_i=0\).
Then \(1+C_{i,\rho}\) is a continuous linear automorphism with
inverse \(1-C_{i,\rho}\), so pullback preserves \(\mathcal S(V)\).

Let \(\psi:k\to\mathbb C^\times\) be a nontrivial additive character,
let \(\xi:k^\times\to\mathbb C^\times\) be any multiplicative
character, and impose the source's exact two laws:
\[
 \begin{aligned}
 l\{\varphi(\gamma Z)\}&=\xi(\det\gamma)l(\varphi)
                      &&(\gamma\in\mathrm{GL}_n(k)),\\
 l\{\varphi(Z+Z_{\cdot,i_0}\rho)\}
       &=\psi(\rho_{i_1})l(\varphi)
                      &&(\rho_{i_0}=0).
 \end{aligned}\tag{LD3}\label{eq:sr-source-local-left-column}
\]
Only Schwartz functions are quantified in these laws.

\begin{srproposition}[The actual two-column support assertion]
\label{prop:sr-source-local-support}
Under (LD3), \(l\) kills the set on which
\(Z_{\cdot,i_0},Z_{\cdot,i_1}\) are linearly independent.
Equivalently it vanishes on every Schwartz test that vanishes
where these columns are dependent. This is \(\mathrm{LocSupp}\).
\end{srproposition}
\begin{proof}
Choose \(t\in k\) with \(\psi(t)\ne1\).
Fix \(Y_0\) with independent columns \(x_0,y_0\).
Choose rows \(\mu,\nu\) with
\[
 \mu x_0=0,\qquad \mu y_0=1,\qquad \nu x_0=1.
\]
They are obtained by extending the two coordinate functionals
from the span of \(x_0,y_0\).
For \(Y\) near \(Y_0\) write \(x=Y_{\cdot,i_0},y=Y_{\cdot,i_1}\), and put
\[
 \begin{split}
 D(Y)&=(\nu x)(\mu y)-(\mu x)(\nu y),\\
 \lambda_Y&=D(Y)^{-1}\{(\nu x)\mu-(\mu x)\nu\},\\
 N_Y&=-t\,x\lambda_Y,\qquad
 \rho_Y(b)=t\,\lambda_Y(Y_{\cdot,b}),\\
 A_Y(Z)&=N_YZ,\qquad B_Y(Z)=Z_{\cdot,i_0}\rho_Y .
 \end{split}\tag{LD4}\label{eq:sr-source-independent-stabilizers}
\]
At \(Y_0\), \(D=1\); thus all these maps are continuous in the
required joint sense on a neighborhood.
Direct multiplication gives
\[
 \lambda_Yx=0,\quad\lambda_Yy=1,\quad
 \rho_Y(i_0)=0,\quad\rho_Y(i_1)=t,\quad
 A_Y^2=B_Y^2=0,\quad A_YB_Y=B_YA_Y.
\]
Moreover \(B_YY=x\rho_Y\), \(A_YY=-x\rho_Y\), and
\(A_YB_YY=0\), so \((1+A_Y)(1+B_Y)Y=Y\).

The matrix \(1+N_Y\) is invertible, with inverse \(1-N_Y\),
and has determinant
\(1-t\lambda_Yx=1\).
The first law in (LD3) therefore gives invariance under \(1+A_Y\).
The second gives multiplication by \(\psi(t)\) under \(1+B_Y\).
Combining the two, with the order in (LD2), gives its scalar
\(c_Y=\psi(t)\ne1\).
Apply Lemma~\ref{lem:sr-source-local-dead} with \(U=\varnothing\).
It shows that \(l\) kills a neighborhood of each independent \(Y_0\).
Lemma~\ref{lem:sr-source-local-kills} proves the assertion on the
whole independent locus.
This is \srcid{kills\_indep}, then \srcid{locSupp}.
\end{proof}

Now fix a third distinct index \(i_2\), and assume in addition
\[
 l\{\varphi(Z+Z_{\cdot,i_1}\rho)\}
       =\psi(\rho_{i_2})l(\varphi)
                \quad(\rho_{i_0}=\rho_{i_1}=0).
 \tag{LD5}\label{eq:sr-source-second-column-law}
\]

\begin{srproposition}[The actual three-column vanishing assertion]
\label{prop:sr-source-local-three}
The laws (LD3)--(LD5) imply \(l(\varphi)=0\) for every
\(\varphi\in\mathcal S(V)\), for every \(n,d\) for which the
three distinct column indices exist.
This is the original \(\mathrm{LocThree}\).
\end{srproposition}
\begin{proof}
First let \(Y_0\) have \(x_0=Y_{0,\cdot,i_0}\ne0\),
and choose a row index \(a_0\) with \(x_0(a_0)\ne0\).
Near this point define, with \(x=Y_{\cdot,i_0},y=Y_{\cdot,i_1}\),
\[
 \rho_Y=e_{i_2}\!\left(-t\,{y(a_0)\over x(a_0)}\right),
 \qquad A_Y=C_{i_1,e_{i_2}t},\qquad B_Y=C_{i_0,\rho_Y}.
 \tag{LD6}\label{eq:sr-source-dependent-stabilizers}
\]
These are jointly continuous near \(Y_0\); they are square zero
and their products in either order are zero.
Let \(U\) be the independent-column locus, already killed by \(l\).
If \(Y\notin U\) near \(Y_0\), then \(x\ne0\) and \(y=sx\)
for one \(s\in k\).
Consequently the changes in column \(i_2\) from \(A_Y\) and \(B_Y\)
cancel, and \((1+A_Y)(1+B_Y)Y=Y\).
Law (LD5) supplies the factor \(\psi(t)\) for \(1+A_Y\).
The second law in (LD3) supplies the factor \(1\) for \(1+B_Y\),
because \(\rho_Y(i_1)=0\).
Thus (LD2) holds with \(c_Y=\psi(t)\ne1\).
Its lemma shows local vanishing near every such \(Y_0\).
The local-kills lemma implies that \(l\) kills the whole locus
\(U_1=\{Y:Y_{\cdot,i_0}\ne0\}\).
This is the source's intermediate \srcid{kills\_ne}.

Finally take any \(Y_0\), and use the fixed operators
\[
                       A_Y=0,\qquad B_Y=C_{i_0,e_{i_1}t}.
\]
They satisfy every continuity and square-zero clause.
Outside \(U_1\), \(Y_{\cdot,i_0}=0\), so
\((1+A_Y)(1+B_Y)Y=Y\).
The second law in (LD3) gives the scalar \(\psi(t)\ne1\).
The moving-stabilizer lemma now gives local vanishing at every
\(Y_0\), hence vanishing everywhere by the local-kills lemma.
This is \srcid{locThree}, with no restriction \(n>0\).
\end{proof}

\paragraph{Frozen-source ledger.}
The compact-coset and moving-stabilizer arguments are
\srcid{src/Development/WeilBound/ThetaLifts/EigenvalueEndgame/CuspidalAverages/ZariskiLocallyAffineLine.lean}
(blob \srcid{c6f5b75d529a8c7e4b5b479465d4dcb5b4fc30af}).
The matrix formulas, independent/dependent support steps and final
vanishing are in
\srcid{src/Development/WeilBound/ThetaLifts/EigenvalueEndgame/CuspidalAverages/LocalFieldT2Space.lean}
(blob \srcid{831e81d9d5ab83074a7c3a318fd6c26e83007527}).
Both were read at the frozen original tag
\srcid{alpoge-source-7cbf7a0c}.
This proves those two local inputs; it does not by itself prove
\(\mathrm{LocTwo}\), \(\mathrm{VacBar}\), the split dictionary,
or any global theta comparison.

\subsection{The original split-place producers of the two cusp keys}
\label{subsec:sr-source-acts-of-reads}

This subsection translates
\srcid{CuspidalAverages.ActsOfReads.keyThree\_of\_locThree} and
\srcid{keyTwo\_of\_locTwo\_vacBar}.
The source inputs \(\mathrm{ThetaBasic},\mathrm{KDecay},\mathrm{IAut}\)
are retained. So are the explicit local distribution statements
\(\mathrm{LocThree},\mathrm{LocTwo}\), the vacuum assertion
\(\mathrm{VacBar}\), and the place assertions \(\mathrm{Places}\).
We write their roles and all constants below.
Proving a cusp key from these inputs does not prove the inputs.

Keep one source triple of first-space ranks \(a,a'\), common target
rank \(d=a+a'\), characters \(\alpha,\beta\), and actual periods
\(I,I'\). Thus a member of the source's comparison space has the
form
\[
 f=I(\Phi)-I'(\Phi'),\qquad
 \Phi\in\mathrm{PolData},\quad
 \Phi'\in\mathrm{PolData}'\cap\mathrm{So}'.
 \tag{AK1}\label{eq:sr-source-key-comparison}
\]
The lower test retains the prescribed \(\mathrm{So}'\) cutoff.

Fix a finite split place \(v\) of \(F\), a map
\(t:E\hookrightarrow F_v=k\), and the source split section
\(\mathrm{sec}:k\times k\to\mathbb A_E\) supported at \(v\).
Its two projections are \(t\) and \(t\circ\sigma\);
these projections determine every element supported at \(v\).
Put
\[
 s_v(x)=\mathrm{sec}(x,0),\qquad
 c_v(h)=t(h)\in\mathrm{GL}_d(k).
\]
Here \(s_v\) includes insertion into the \(v\)-component.
The section also gives the source local character
\(\chi_v=\chi\circ s_v\).
If \(\chi\) satisfies \(\mathrm{IsChar}\), then
\(\chi_v\ne1\): the nontrivial element supported at \(v\),
with second projection zero, is \(s_v\) of its first projection.
This proves \srcid{psi\_ne}; it does not use only rational
modulations, on which a global additive character would vanish.

\paragraph{The precise local period maps.}
Write \(\mathcal S(M_{n,d}(k))\) for locally constant compactly
supported functions on \(n\)-by-\(d\) matrices.
For cofinally all finite \(v\), the source constructs maps
\(S:\mathcal S(M_{a,d}(k))\to\{\text{automorphic functions on }U(W)\}\)
and \(S'\) with \(a'\) in place of \(a\), together with local
model changes \(\iota,\iota'\), such that
\[
 S(\iota\,\mathrm{vac}_v)=I(\Phi),\qquad
 S'(\iota'\,\mathrm{vac}'_v)=I'(\Phi').
 \tag{AK2}\label{eq:sr-source-local-period-values}
\]
The maps are additive on Schwartz functions and homogeneous.
For the large map their exact transformation laws are
\[
 \begin{split}
 S\{\varphi(\gamma Y)\}
       &=\xi(\det\gamma)\,S\varphi,\\
 R_h(S\varphi)&=\eta(\det c_v(h))\,
                     S\{\varphi(Yc_v(h))\}
                   \quad(h\text{ supported at }v),\\
 \xi&=\{\alpha_v^{-1}\beta_v|\cdot|_v^{d/2}\}^{-1},
 \qquad \eta=\beta_v|\cdot|_v^{a/2}.
 \end{split} \tag{AK3}\label{eq:sr-source-local-period-laws}
\]
For the small map replace
\[
 (\xi,\eta,a)=
 \bigl(\{\alpha_v\beta_v^{-1}|\cdot|_v^{d/2}\}^{-1},
              \alpha_v|\cdot|_v^{a'/2},a'\bigr).
 \tag{AK4}\label{eq:sr-source-small-local-period-laws}
\]
We view the unit-valued characters through their complex values.
The absolute value has \(|\varpi|_v=q^{-1}\).

Here is the construction in \srcid{exists\_loc},
\srcid{locV}, and \srcid{locV'}.
For a global test \(\Phi\), define insertion by
\[
 \mathrm{ins}_v(\Phi,\varphi)(x)
   =\varphi(x_v)\,\Phi(x^{v,0}),
\]
where \(x^{v,0}\) replaces its \(v\)-coordinate by zero.
For a pure restricted tensor, insertion changes exactly one factor.
Outside a finite set that factor equals the vacuum and has value
one at zero. Hence insertion of the vacuum recovers \(\Phi\).
Insertion of any local Schwartz function remains in
\(\mathrm{PolData}\): its archimedean polynomial factor is unchanged,
and all finite factors remain Schwartz and vacuum almost everywhere.
The set of tests with these two properties is a vector subspace,
so the assertion extends to every finite sum of pure tensors.
For the source's \(\mathrm{So}\), exclude also its one distinguished
finite place \(v_1\). The local condition at \(v_1\) is then unchanged.
Intersecting the two insertion-stable subspaces proves the assertion
for \(\mathrm{PolData}'\cap\mathrm{So}'\).

The actual local Weil dictionary supplies \(\iota\), its inverse
\(j\) on Schwartz functions, and the intertwining laws.
Define
\[
                  S\varphi=I(\mathrm{ins}_v(\Phi,j\varphi)).
 \tag{AK5}\label{eq:sr-source-local-period-construction}
\]
The insertion argument puts every test used here in the source
period domain. \(\mathrm{KDecay}\) gives convergence; consequently
theta-sum linearity and linearity of the convergent integral give
additivity. \(\mathrm{IAut}\) supplies automorphy.
Intertwining the first-group action and applying invariance of its
period gives the first law in (AK3). The local transpose used there
has the same determinant as \(\gamma\); the inverse character is
therefore exactly the displayed \(\xi\).
Intertwining the second-group action and right translation gives
the second law and its displayed \(\eta\).
The identical calculation with the small source pair gives (AK4).
No scalar in the local dictionary is independently renormalized here.
The dictionary's model change and inverse, including its splitting
characters, are retained construction inputs in this step.

\paragraph{From a global Whittaker functional to a local distribution.}
Let \(u,u'\) be a source frame of size \(j\):
\[
 u^tJ\bar u=0,\qquad (u')^tJ\bar u'=0,\qquad u^tJ\bar u'=1.
\]
There are an injection \(e:\{0,\ldots,j-1\}\to\{0,\ldots,d-1\}\)
and \(c\in\mathrm{GL}_d(E)\) such that
\[
 c_{\cdot,e(i)}=u_{\cdot,i},\qquad
 (J\bar u')^t c=(1_d)_{e(\cdot),\cdot}.
 \tag{AK6}\label{eq:sr-source-frame-complement}
\]
Indeed the displayed left inverse of \(u\) splits \(E^d\) into
its image and the kernel of that left inverse. Append a basis of
the kernel to the specified columns of \(u\).
This proves \srcid{frame\_compl}.

Let \(\Lambda\) satisfy the actual \(\mathrm{IsWh}\) laws for this
frame and \(\chi\). Fix any \(g\in U(W)(\mathbb A)\).
Every local matrix \(c_v(g)\) has a lift \(g_v\) supported at \(v\).
It is obtained by putting that matrix in the first split factor
and its inverse \(J\)-adjoint in the conjugate split factor.
Then
\[
                       g=r g_v,
\]
where \(r\) commutes with all elements supported at \(v\):
both split components of \(r\) are \(1\).
The equality follows from the two-projection uniqueness of the split
section and the unitary equation. This is \srcid{ex\_split}.

On the Schwartz subspace define
\[
 l_0(\varphi)=
       \Lambda\bigl(R_r\,S\{\varphi(Yc)\}\bigr).
\]
Automorphy of \(S\varphi\), and additivity of \(\Lambda\) on
automorphic functions, make this linear.
Extend it algebraically to a linear functional \(l\) on all
matrix functions; only its restriction to Schwartz functions is used.
The original \srcid{exists\_l} proves
\[
 \begin{split}
 l\{\varphi(\gamma Y)\}&=\xi(\det\gamma)l(\varphi),\\
 l\{\varphi(\mathrm{colOp}_{e(i),\rho}Y)\}
       &=\chi_v(\rho_{e(i+1)})l(\varphi)
       \quad\bigl(\rho_{e(b)}=0\text{ for }b\le i\bigr),\\
 \Lambda(R_{kg}S\varphi)
       &=\eta(\det(c_v(k)c_v(g)))\,
          l\{\varphi(Yc^{-1}c_v(k)c_v(g))\}
                  \quad(k\text{ supported at }v),
 \end{split}\tag{AK7}\label{eq:sr-source-local-whittaker-distribution}
\]
where
\[
 \mathrm{colOp}_{i,\rho}Y
                  =Y+Y_{\cdot,i}\rho .
\]
In (AK7), \(c\) denotes its image under \(t\).

For completeness, the column law uses the local matrix
\(N=1+e_{e(i)}\rho\). Its determinant is \(1+\rho_{e(i)}=1\).
Lift \(cNc^{-1}\) to the split unitary group at \(v\).
The frame columns transform by an upper unipotent matrix whose
consecutive superdiagonal sum is \(s_v(\rho_{e(i+1)})\);
the earlier consecutive entries vanish because
\(\rho_{e(b)}=0\) for \(b\le i\).
On the orthogonal complement of the frame span, the displacement
is a multiple of \(u_{\cdot,i}\), hence belongs to that span.
These are all clauses of \(\mathrm{IsWh.law}\).
Applying that law, then the second law in (AK3), proves the column
identity. The determinant factor is \(1\).
The first line of (AK7) follows from the first line of (AK3).
For the last line substitute \(g=r g_v\), commute \(r\) past \(k\),
and use the second line of (AK3). Thus no continuity of \(l\) is
assumed or required.

\begin{srproposition}[The source's three-column key producer]
\label{prop:sr-source-key-three-producer}
Assume \(\mathrm{ThetaBasic},\mathrm{KDecay},\mathrm{IAut},
\mathrm{LocThree},\mathrm{Places}\).
Then \(\mathrm{KeyThree}\) holds with its full quantifiers:
every frame of size at least three, every admissible
\(\chi,\Lambda\), every \(f\) in (AK1), and every right translate.
\end{srproposition}
\begin{proof}
There are infinitely many split places by \(\mathrm{Places}\).
Choose one outside the two finite exceptional sets in (AK2).
The functional \(l\) in (AK7) has the left covariance and the two
successive column laws required by \(\mathrm{LocThree}\), at the
distinct indices \(e(0),e(1),e(2)\).
The local character is nontrivial, as proved above.
Thus \(\mathrm{LocThree}\) makes \(l\) vanish on every Schwartz test.
Take \(k=1\) in (AK7). Its right matrix substitution preserves the
Schwartz space, so \(\Lambda(R_gS\varphi)=0\).
Apply this separately to \(S(\iota\mathrm{vac})\) and
\(S'(\iota'\mathrm{vac}')\).
Both functions are automorphic. Additivity and homogeneity of
\(\Lambda\) therefore give
\(\Lambda(R_g(I\Phi-I'\Phi'))=0\).
This is exactly \srcid{keyThree\_of\_locThree}.
\end{proof}

\paragraph{The numerical sequence used for the two-column key.}
For \(q>1,\ r\in\mathbb C^\times,\ c\in\mathbb Z\), put
\[
 Y_m(q,r,c)=(qr)^{-m}
       \sum_{0\le i<(2m-c+1)_+}r^i,\qquad
                     (x)_+=\max\{x,0\}.
 \tag{AK8}\label{eq:sr-source-key-two-sequence}
\]
Integer powers and the truncation in this formula are literal.

\begin{srlemma}\label{lem:sr-source-key-two-sequence}
For every \(m\ge0\),
\[
 Y_m(q,r^{-1},c)=r^cY_m(q,r,c).
 \tag{AK9}\label{eq:sr-source-key-two-inversion}
\]
If \(|r|>q\), then \(|Y_m(q,r,c)|\to\infty\).
\end{srlemma}
\begin{proof}
When \(2m-c+1\le0\), both sums are zero.
Otherwise write \(2m-c=N\ge0\).
Reversing the finite sum gives
\(\sum_{i=0}^Nr^{-i}=r^{-N}\sum_{i=0}^Nr^i\).
Since \(c=2m-N\), multiplication by the two prefactors gives (AK9).

Put \(R=|r|>q>1\).
For all sufficiently large \(m\), \(N=2m-c+1\ge0\), and
the finite geometric identity and the norm triangle inequality give
\[
 R^N-1
   \le \left|\sum_{i=0}^{N-1}r^i\right|(R+1).
\]
Multiplying by \((qR)^{-m}\), and using
\((qR)^{-m}\le1\), yields
\[
 |Y_m(q,r,c)|
       \ge {R^{1-c}(R/q)^m-1\over R+1}.
\]
The right side tends to infinity.
These are the source's \srcid{yy\_inv} and
\srcid{yy\_tendsto}; no phase convergence is asserted.
\end{proof}

\begin{srproposition}[The source's two-column key producer]
\label{prop:sr-source-key-two-producer}
Assume \(\mathrm{ThetaBasic},\mathrm{KDecay},\mathrm{IAut},
\mathrm{LocTwo},\mathrm{VacBar},\mathrm{Places}\).
For every size-two source frame, every specified Levi map \(m\),
every admissible \(\chi,\Lambda\), every comparison form \(f\),
and every \(h_0\in U(W)(\mathbb A)\), there are a neighborhood
\(V\) of \(h_0\), matrices \(b_i\in\mathrm{GL}_2(\mathbb A_E)\),
and complex numbers \(y_i\), with
\[
 b_0=1,\quad \det b_i=1,\quad |y_i|\longrightarrow\infty,
 \qquad
 \forall h\in V\ \exists x_h\in\mathbb C\
 \forall i\ge0:\quad
              \Lambda(R_{m(b_i)h}f)=x_hy_i.
 \tag{AK10}\label{eq:sr-source-key-two-conclusion}
\]
\end{srproposition}
\begin{proof}
Choose the rational completion \(c\) in (AK6) first.
The place assertions permit choosing one split \(v\) for which:
both (AK2) hold; both model changes carry the vacuum to
\(\kappa1_{M_{a,d}(\mathcal O_v)}\) and
\(\kappa'1_{M_{a',d}(\mathcal O_v)}\);
the entries of \(c,c^{-1},c_v(h_0),c_v(h_0^{-1})\) are integral;
and \(\alpha_v,\beta_v\) are unramified.
Indeed all these requirements exclude only finitely many places,
whereas the split places are infinite.
The two vacuum scalars are not assumed nonzero.

Choose a uniformizer \(\varpi\), let \(q=|\mathcal O_v/\varpi|\),
and choose the exact conductor \(c_0\in\mathbb Z\) of
\(\chi_v\): it is trivial on \(\varpi^{c_0}\mathcal O_v\) and
nontrivial on \(\varpi^{c_0-1}\mathcal O_v\).
Here is the source's short proof of \srcid{exists\_cond}.
Continuity at zero gives an ideal \(\varpi^m\mathcal O_v\)
on which \(|\chi_v(x)-1|<1/2\).
Every nonnegative integer multiple of a point of that ideal
remains in the ideal. Thus, for \(z=\chi_v(x)\),
\(|z^n-1|<1/2\) for every \(n\ge0\).
The identity
\[
 z^{2^{j+1}}-1=(z^{2^j}-1)\{2+(z^{2^j}-1)\}
\]
implies
\[
 |z^{2^j}-1|\ge(3/2)^j|z-1|.
\]
The left side is bounded by \(1/2\), so \(z=1\).
There is therefore at least one exponent on whose ideal
the character is trivial. Fixing one \(x\) with
\(\chi_v(x)\ne1\) bounds these exponents below, since
\(x\in\varpi^c\mathcal O_v\) for all sufficiently negative
integers \(c\).
Their least integer \(c_0\) has the asserted triviality;
minimality gives a point of
\(\varpi^{c_0-1}\mathcal O_v\) with nontrivial character.

Use the split section to choose \(b_i\) whose first local projection
is \(\operatorname{diag}(\varpi^i,\varpi^{-i})\), whose conjugate
projection is \(1\), and which is \(1\) away from \(v\).
Then \(b_0=1\), and \(\det b_i=1\).
For the prescribed Levi map the frame formula gives
\[
 c_v(m(b_i))c
   =c\,\mathrm{torus}_{e(0),e(1)}(\varpi^i).
 \tag{AK11}\label{eq:sr-source-key-two-torus}
\]
The source proves this by substituting the two projections into
\(\mathrm{levM}\); (AK6) cancels the frame row and column matrices.
This is \srcid{exists\_torus}; the matrices \(b_i\) are the same for
both source periods.

Set
\[
 V=\{h:c_v(h),\,c_v(h^{-1})\in M_d(\mathcal O_v)\}.
\]
The valuation ring is open, and projection, inversion and matrix
entries are continuous. Hence \(V\) is open and contains \(h_0\).
For \(h\in V\), apply (AK7) with \(g=h\).
Integral invertibility of \(c\) and \(c_v(h)\) gives
\[
 1_{\mathrm{int}}\{Yc^{-1}c_v(m(b_i))c_v(h)\}
       =1_{\mathrm{int}}\{
                    Y\,\mathrm{torus}_{e(0),e(1)}(\varpi^i)\}.
 \tag{AK12}\label{eq:sr-source-key-two-vacuum-torus}
\]
To verify this, factor the matrix on the left as
\(\mathrm{torus}_{e(0),e(1)}(\varpi^i)c^{-1}c_v(h)\)
using (AK11).
Right multiplication by the final matrix preserves integral matrices
in both directions. This is \srcid{oneO\_tor}.
The determinant of \(c_v(m(b_i))\) is one, so its \(\eta\)-factor
in (AK7) has no \(i\)-dependence.
The remaining \(\eta(\det c_v(h))\) and vacuum scalar are absorbed
into a constant depending on \(h\).

Write \(a=a'+2\epsilon\), using the triple's parity and \(a'<a\), and
put
\[
 r={\alpha_v(\varpi)\over\beta_v(\varpi)}q^\epsilon.
 \tag{AK13}\label{eq:sr-source-key-two-r}
\]
The triple's rank gap gives \(\epsilon>1\).
Its Hecke character data and the split section give
\(|\alpha_v(\varpi)|=|\beta_v(\varpi)|=1\).
Thus \(|r|=q^\epsilon>q\).
All these statements are made after fixing the actual triple;
no bound uniform in a changing triple is inferred.

The large left character in (AK3), at \(\varpi\), is
\((\alpha_v/\beta_v)(\varpi)q^{d/2}\).
Consequently its local \(\mathrm{LocTwo}\) parameter is
\[
                    \xi(\varpi)q^{a-d}
                       =(\alpha_v/\beta_v)(\varpi)q^\epsilon=r.
\]
For (AK4) the parameter is
\(\xi'(\varpi)q^{a'-d}=r^{-1}\).
The characters \(\xi,\xi'\) are unramified: both original
characters and the absolute-value powers are trivial on local units.
Thus every premise of the local \(\mathrm{LocTwo}\) formula is present.
Applied to (AK12), that formula gives
\[
 \Lambda(R_{m(b_i)h}I\Phi)=A_hY_i(q,r,c_0),\qquad
 \Lambda(R_{m(b_i)h}I'\Phi')=B_hY_i(q,r^{-1},c_0).
\]
Their local additive characters and conductor are identical.
Apply (AK9), and subtract, using automorphy for the additivity of
\(\Lambda\).
With
\[
             y_i=Y_i(q,r,c_0),\qquad
                         x_h=A_h-B_hr^{c_0},
\]
we obtain (AK10).
The preceding lemma proves \(|y_i|\to\infty\).
This is the full quantifier order of
\srcid{keyTwo\_of\_locTwo\_vacBar}: \(V,b_i,y_i\) precede \(h\),
whereas \(x_h\) may depend on \(h\).
\end{proof}

\paragraph{Frozen source and retained leaves.}
All source files below have prefix
\srcid{src/Development/WeilBound/}.
The local-to-global construction is
\srcid{ThetaLifts/EigenvalueEndgame/CuspidalAverages/ActsOfReads.lean}
(blob \srcid{4e02838d8c6835af8c49065256decc48a0b20f9c});
its conductor and sequence calculations are
\srcid{ThetaLifts/EigenvalueEndgame/CuspidalAverages/ExistsCond.lean}
(blob \srcid{123a7a68b9cdf0725772ddb17331e96dc4a2d617}).
Both were read at \srcid{alpoge-source-7cbf7a0c}.

The still separate local leaves are
\srcid{LocalFieldT2Space.locSupp/locThree},
\srcid{ColOpSingle.locTwo\_of}, and
\srcid{VacIntZ.vacBar}, in that same cusp directory.
The place assertions and the exact split Weil dictionary, including
\srcid{dictionary\_extended} and \srcid{section\_statement},
also remain separate producers.
This subsection proves their application to the actual large/small
periods, and the shared sequence and neighborhood; it does not replace
one of those producers by a choice of normalizations.
